% ============================================================
\section{Ejercicios de la Unidad 1}
% ============================================================

Los siguientes ejercicios están organizados por tema y nivel de dificultad. Se recomienda resolverlos sin apoyo de software en una primera etapa y utilizar herramientas computacionales únicamente para verificación.

\subsection{Potencias de la unidad imaginaria}
\begin{enumerate}
\item Calcule $i^{17}$.
\item Calcule $i^{38}$.
\item Calcule $i^{73}$.
\item Calcule $i^{2026}$.
\item Simplifique $i^{11}+i^{22}+i^{33}+i^{44}$.
\item Simplifique $i^{15}-2i^{26}+3i^{37}-i^{48}$.
\end{enumerate}

\subsection{Forma binómica, parte real e imaginaria}
\begin{enumerate}
\item Determine $\operatorname{Re}(z)$ y $\operatorname{Im}(z)$ para $z=5-7i$.
\item Determine $\operatorname{Re}(z)$ y $\operatorname{Im}(z)$ para $z=-3+\frac{5}{2}i$.
\item Escriba $-8$ en forma binómica e identifique sus partes real e imaginaria.
\item Escriba $6i$ en forma binómica e identifique sus partes real e imaginaria.
\item Clasifique como real, imaginario puro o complejo no real: $4$, $-3i$, $2+7i$, $0$, $5-2i$.
\item Determine $a,b\in\mathbb{R}$ si $z=(2a-1)+(3b+4)i=7-5i$.
\end{enumerate}

\subsection{Igualdad de números complejos}
\begin{enumerate}
\item Determine $x$ y $y$ si $(x+2)+(y-1)i=5+4i$.
\item Determine $x$ y $y$ si $(2x-3)+(3y+2)i=7-10i$.
\item Determine $a$ y $b$ si $(a-b)+(a+b)i=4+2i$.
\item Determine $x$ y $y$ si $(x^2-1)+(2y+3)i=8+7i$ y $x>0$.
\item Determine $x$ y $y$ si $(3x+y)+(x-2y)i=7-i$.
\item Pruebe que si $a+bi=0$, entonces $a=b=0$.
\end{enumerate}

\subsection{Suma y resta de números complejos}
\begin{enumerate}
\item Calcule $(3+5i)+(7-2i)$.
\item Calcule $(-4+9i)+(6-11i)$.
\item Calcule $(8-3i)-(2+7i)$.
\item Simplifique $(2+i)-(5-3i)+(7+4i)$.
\item Sean $z_1=3-2i$, $z_2=-1+5i$, $z_3=4+i$. Calcule $2z_1-z_2+3z_3$.
\item Verifique explícitamente la asociatividad de la suma para $z_1=1+i$, $z_2=2-3i$, $z_3=-4+2i$.
\end{enumerate}

\subsection{Producto de números complejos}
\begin{enumerate}
\item Calcule $(2+3i)(4-i)$.
\item Calcule $(1-5i)(-2+3i)$.
\item Calcule $(3+2i)^2$.
\item Simplifique $(1+i)^4$ usando primero multiplicación directa.
\item Calcule $(2-i)(2+i)(1+3i)$.
\item Determine todos los $z=a+bi$ tales que $(1+i)z=4+2i$.
\end{enumerate}

\subsection{Conjugado complejo}
\begin{enumerate}
\item Determine $\overline z$ para $z=7-5i$.
\item Calcule $z+\overline z$ y $z-\overline z$ para $z=-3+4i$.
\item Exprese $\operatorname{Re}(z)$ en términos de $z$ y $\overline z$.
\item Exprese $\operatorname{Im}(z)$ en términos de $z$ y $\overline z$.
\item Verifique que $\overline{zw}=\overline z\,\overline w$ para $z=2+i$, $w=3-4i$.
\item Demuestre que $z$ es real si y sólo si $z=\overline z$.
\end{enumerate}

\subsection{Módulo de números complejos}
\begin{enumerate}
\item Calcule $|3+4i|$.
\item Calcule $|-5+12i|$.
\item Calcule $|1-\sqrt3\,i|$.
\item Encuentre todos los $z=a+bi$ tales que $|z|=5$ y $a=3$.
\item Demuestre que $|z|^2=z\overline z$.
\item Demuestre que $|z|=0$ si y sólo si $z=0$.
\end{enumerate}

\subsection{División e inverso multiplicativo}
\begin{enumerate}
\item Calcule $\dfrac{1}{2+i}$.
\item Calcule $\dfrac{3+2i}{1-i}$.
\item Calcule $\dfrac{4-5i}{2+3i}$.
\item Determine $(3-4i)^{-1}$.
\item Simplifique $\dfrac{1+i}{1-i}$.
\item Demuestre que para $z\neq0$, $z^{-1}=\dfrac{\overline z}{|z|^2}$.
\end{enumerate}

\subsection{El plano complejo y lugares geométricos}
\begin{enumerate}
\item Represente en el plano complejo $z=3+2i$, $z=-2+4i$, $z=-3i$ y $z=5$.
\item Describa geométricamente el conjunto $\{z\in\mathbb{C}: |z|=4\}$.
\item Describa geométricamente $\{z: |z-(2-i)|<3\}$.
\item Describa geométricamente $\{z: |z+1+2i|\le 2\}$.
\item Determine el lugar geométrico de los puntos que satisfacen $|z-1|=|z+1|$.
\item Determine el lugar geométrico de los puntos que satisfacen $|z-i|=2|z+i|$.
\end{enumerate}

\subsection{Distancia y desigualdades}
\begin{enumerate}
\item Calcule la distancia entre $z_1=1+i$ y $z_2=4+5i$.
\item Calcule la distancia entre $z_1=-2+3i$ y $z_2=1-i$.
\item Verifique la desigualdad triangular para $z=1+i$ y $w=1-i$.
\item Verifique la desigualdad triangular inversa para $z=3+4i$ y $w=1$.
\item Demuestre que $|z+w|\le |z|+|w|$.
\item Demuestre la identidad del paralelogramo: $|z+w|^2+|z-w|^2=2(|z|^2+|w|^2)$.
\end{enumerate}

\subsection{Argumento y argumento principal}
\begin{enumerate}
\item Determine $\arg z$ y $\operatorname{Arg}(z)$ para $z=1+i$.
\item Determine $\arg z$ y $\operatorname{Arg}(z)$ para $z=-1+i$.
\item Determine $\arg z$ y $\operatorname{Arg}(z)$ para $z=-1-i$.
\item Determine $\arg z$ y $\operatorname{Arg}(z)$ para $z=-4$.
\item Determine $\arg z$ y $\operatorname{Arg}(z)$ para $z=-3i$.
\item Explique por qué $\arg z$ es multivaluado y $\operatorname{Arg}(z)$ no.
\end{enumerate}

\subsection{Conversión entre forma rectangular y polar}
\begin{enumerate}
\item Convierta $1+i$ a forma polar.
\item Convierta $-1+\sqrt3\,i$ a forma polar.
\item Convierta $-2-2i$ a forma polar.
\item Convierta $4\left(\cos\frac{\pi}{3}+i\sin\frac{\pi}{3}\right)$ a forma rectangular.
\item Convierta $5\left(\cos\frac{5\pi}{6}+i\sin\frac{5\pi}{6}\right)$ a forma rectangular.
\item Exprese $-3i$ en forma polar utilizando un argumento principal y un argumento general.
\end{enumerate}

\subsection{Producto y cociente en forma polar}
\begin{enumerate}
\item Multiplique $2\operatorname{cis}\frac{\pi}{6}$ y $3\operatorname{cis}\frac{\pi}{4}$.
\item Divida $6\operatorname{cis}\frac{5\pi}{6}$ entre $3\operatorname{cis}\frac{\pi}{3}$.
\item Calcule el módulo y argumento de $(1+i)(\sqrt3+i)$ sin convertir primero a forma rectangular.
\item Calcule $\dfrac{2e^{i\pi/3}}{4e^{-i\pi/6}}$.
\item Si $|z_1|=3$, $|z_2|=5$, $\arg z_1=\frac{\pi}{7}$ y $\arg z_2=-\frac{\pi}{4}$, determine módulo y argumento de $z_1z_2$.
\item Demuestre en forma polar que $|z_1z_2|=|z_1||z_2|$.
\end{enumerate}

\subsection{Fórmula de Euler y forma exponencial}
\begin{enumerate}
\item Evalúe $e^{i\pi/6}$.
\item Evalúe $e^{i3\pi/4}$.
\item Exprese $1-i$ en forma exponencial.
\item Exprese $-2+2i$ en forma exponencial.
\item Demuestre que $\cos\theta=\dfrac{e^{i\theta}+e^{-i\theta}}{2}$.
\item Demuestre que $\sin\theta=\dfrac{e^{i\theta}-e^{-i\theta}}{2i}$.
\end{enumerate}

\subsection{Teorema de De Moivre y potencias}
\begin{enumerate}
\item Calcule $(1+i)^8$ usando De Moivre.
\item Calcule $(-1+i)^6$ usando De Moivre.
\item Calcule $(\sqrt3+i)^{10}$.
\item Calcule $\left(2e^{i\pi/5}\right)^7$.
\item Use De Moivre para deducir una fórmula para $\cos(3\theta)$.
\item Use De Moivre para deducir una fórmula para $\sin(4\theta)$.
\end{enumerate}

\subsection{Raíces $n$-ésimas}
\begin{enumerate}
\item Encuentre las raíces cuadradas de $i$.
\item Encuentre las raíces cúbicas de $8$.
\item Encuentre las raíces cúbicas de $8i$.
\item Encuentre las raíces cuartas de $-16$.
\item Encuentre las cinco raíces de $32$.
\item Demuestre que las raíces $n$-ésimas de un número complejo no nulo están igualmente espaciadas angularmente.
\end{enumerate}

\subsection{Raíces de la unidad}
\begin{enumerate}
\item Encuentre las raíces cúbicas de la unidad y represéntelas geométricamente.
\item Encuentre las raíces cuartas de la unidad.
\item Encuentre las raíces sextas de la unidad.
\item Determine cuáles de las raíces octavas de la unidad son primitivas.
\item Demuestre que la suma de las raíces $n$-ésimas de la unidad es cero para $n>1$.
\item Factorice $z^5-1$ completamente sobre $\mathbb{C}$.
\end{enumerate}

\subsection{Ecuaciones polinómicas}
\begin{enumerate}
\item Resuelva $z^2+1=0$.
\item Resuelva $z^2-2z+5=0$.
\item Resuelva $z^3-8=0$.
\item Resuelva $z^4+16=0$.
\item Factorice $z^4-1$ sobre $\mathbb{C}$.
\item Construya un polinomio mónico real de grado $4$ cuyas raíces incluyan $1+2i$ y $-3+i$.
\end{enumerate}

\subsection{Teorema Fundamental del Álgebra y conjugación}
\begin{enumerate}
\item Explique qué significa que $\mathbb{C}$ sea algebraicamente cerrado.
\item ¿Cuántas raíces posee, contando multiplicidad, un polinomio complejo de grado $7$?
\item Si $2+3i$ es raíz de un polinomio real, ¿qué otra raíz debe aparecer?
\item Construya el factor cuadrático real asociado al par de raíces $4\pm5i$.
\item Demuestre que las raíces no reales de un polinomio real aparecen en pares conjugados.
\item Explique por qué todo polinomio real de grado impar tiene al menos una raíz real.
\end{enumerate}

\subsection{Exponencial y logaritmo complejo}
\begin{enumerate}
\item Calcule $e^{1+i\pi}$.
\item Calcule $e^{2+i\pi/2}$.
\item Resuelva $e^z=1$.
\item Determine todos los valores de $\log(-1)$.
\item Calcule $\operatorname{Log}(1+i)$.
\item Calcule $\operatorname{Log}(-2i)$.
\end{enumerate}

\subsection{Problemas integradores}
\begin{enumerate}
\item Sea $z=1+i$. Calcule sucesivamente $\overline z$, $|z|$, $\operatorname{Arg}(z)$, su forma polar, su forma exponencial, $z^8$ y sus raíces cúbicas.
\item Sea $z=-1+\sqrt3\,i$. Exprese $z$ en forma polar y exponencial y calcule $z^{12}$.
\item Resuelva completamente $z^4=16i$ y represente geométricamente las soluciones.
\item Determine el lugar geométrico definido por $|z-1|+|z+1|=4$ e interprete el resultado.
\item Sea $p(z)=z^4+5z^2+4$. Factorice completamente $p$ sobre $\mathbb{R}$ y sobre $\mathbb{C}$.
\item Pruebe que la transformación $T(z)=e^{i\theta}z$ preserva distancias: $|T(z_1)-T(z_2)|=|z_1-z_2|$.
\end{enumerate}

\subsection{Ejercicios tipo examen}
\begin{enumerate}
\item Sin usar calculadora, simplifique $\dfrac{(1+i)^{10}}{(1-i)^6}$.
\item Determine todas las soluciones de $z^3=-8i$ y expréselas en forma rectangular.
\item Sea $z$ tal que $|z|=2$ y $\operatorname{Arg}(z)=\frac{2\pi}{3}$. Calcule $z^5$.
\item Encuentre todos los $z$ que satisfacen simultáneamente $|z|=2$ y $|z-2|=2$.
\item Factorice completamente $z^6-1$ sobre $\mathbb{C}$.
\item Determine todos los valores de $w$ tales que $e^w=1+i$.
\end{enumerate}

% ============================================================
\subsection{Sugerencia de uso docente}
% ============================================================

\begin{itemize}
\item \textbf{Nivel básico:} seleccionar ejercicios de cálculo directo al inicio de cada subtema.
\item \textbf{Nivel intermedio:} combinar representación algebraica, geométrica y polar.
\item \textbf{Nivel avanzado:} utilizar demostraciones, lugares geométricos, raíces y problemas integradores.
\item \textbf{Evaluación:} los ejercicios tipo examen están formulados sin indicar explícitamente el método a emplear.
\end{itemize}